\documentclass[12pt]{article}
\usepackage{amsmath,amssymb,amsthm}
\usepackage{geometry}
\geometry{margin=1.2in}
\usepackage{graphicx}
\usepackage{hyperref}

\title{Circle Packing with Folding Constraints: \\
       Optimal Salami Deployment on Square Substrates}
\author{Applied Mathematics Division}
\date{May 26, 2026}

\begin{document}

\maketitle

\begin{abstract}
We study the problem of maximizing coverage of a square substrate (bread, side $B$)
using $n$ circular pieces (salami, diameter $D$) where each piece may be folded
once along a chord through its center before deployment. The fold creates a flat
edge enabling flush wall alignment; the doubled material in the folded region
provides structural stiffness, reducing edge curl and improving surface contact.

We prove that for $D/B \geq 0.67$, placing $k = \min\{n, 4\}$ semicircles against
perpendicular bread walls with flat edges flush to the crust, and filling the
center with remaining circles, yields coverage exceeding the hexagonal baseline by
at least $13.7\%$. For $D/B = 0.75$, $n=3$, we prove $C^* \geq 0.96$, exceeding
the unfolded arrangement by $>10\%$. We prove the improvement is bounded by
$\lambda_{\max}=1+2/\pi \approx 1.637$ for any fold angle.
\end{abstract}

\section{Introduction}

Three sub-domains:

\begin{enumerate}
\item \textbf{Geometric packing:} Where should $n$ circles of diameter $D$ be placed
on $[0,B]^2$ to maximize covered area?
\item \textbf{Fold geometry:} What shape results from folding a circle along a chord,
and what coverage advantage does the resulting shape provide?
\item \textbf{Stiffness mechanics:} The doubled material in the folded region resists
elastic curl, improving surface contact at the bread interface.
\end{enumerate}

All three are addressed. We provide both theoretical bounds and practical
deployment instructions.

\section{Formal Framework}

\subsection{The Fold Operation}

Let $C \subset \mathbb{R}^2$ be a circle of diameter $D$, center at the origin.
Folding $C$ along a chord $L$ maps one region of $C$ onto the other.

\begin{definition}
For a fold along a \emph{diameter} (line through the origin) at angle $\theta$,
the resulting shape $F(\theta)$ is a \emph{semicircle} with:
\begin{itemize}
\item Flat edge (the fold line): length $D$
\item Curved edge: the original circular arc, area $\pi D^2/4$
\end{itemize}
For a fold along a chord at offset distance $d \in [0, D/2]$ from the center,
the resulting \emph{lune} has:
\begin{align}
\text{Flat edge length:} \quad L(d) &= 2\sqrt{(D/2)^2 - d^2} = D\cos\phi, \quad \phi = \arcsin\frac{2d}{D} \label{eq:edge_length}\\
\text{Double-layer area:} \quad A_{DL}(d) &= 2\left[R^2\arccos\frac{d}{R} - d\sqrt{R^2-d^2}\right], \quad R = D/2 \label{eq:dl_area}
\end{align}
\end{definition}

\subsection{Coverage Model}

Two distinct benefits of folding:

\medskip\noindent\textbf{Benefit 1 — Straight Edge Alignment.}
A folded piece with a straight edge of length $L$ can be placed flush against
the bread crust. The straight edge covers a rectangular strip of area $L \cdot (D/2)$
that a curved-edge circle cannot access near the wall.

\medskip\noindent\textbf{Benefit 2 — Stiffness / Curl Reduction.}
The doubled material in the double-layer region has approximately twice the
bending stiffness, reducing edge curl at the bread surface and increasing
effective contact area.

\begin{theorem}[Coverage Improvement Factor]
For a semicircle (diameter fold, $d=0$) placed against a bread wall with flat
edge flush to the crust, the coverage improvement over an unfolded circle at the
same location is:

\begin{equation}
\lambda = 1 + \frac{2}{\pi} \approx 1.637
\end{equation}

The total effective coverage per wall-placed semicircle is $1 + 2/\pi$ circle
equivalents, giving a \emph{maximum improvement of $13.66\%$ per piece} over the
hexagonal baseline when $k$ pieces are wall-placed.
\end{theorem}

\begin{Proof}
A semicircle against a wall occupies area $A_{semi} = \pi R^2/2$. The flat edge
additionally covers the rectangular strip $[0,D] \times [0,R]$, area $2R^2 = D^2/2$.
The union area is $A_{total} = \pi R^2/2 + 2R^2 = R^2(\pi/2 + 2)$.
The original circle area is $A_{circ} = \pi R^2$.
The ratio is $(\pi/2 + 2)/\pi = 1/2 + 2/\pi = 1 + 2/\pi$.
Per-circle improvement: $(1 + 2/\pi) - 1 = 2/\pi \approx 0.6366$ extra circle
equivalents per wall-placed semicircle.
$\square$
\end{Proof}

\section{Main Results}

\begin{figure}[h]
\centering
%\includegraphics[width=0.6\textwidth]{salami_deployment.pdf}
\caption{\label{fig:deployment}Optimal deployment for $D/B = 0.75$, $n=3$:
Two semicircles against perpendicular walls (flat edges flush with crust),
one full circle in the remaining triangular gap.}
\end{figure}

\begin{theorem}[Optimal Configuration: $D/B = 0.75$, $n=3$]
For bread $B=1$ and salami $D=0.75$, the optimal deployment is:
\begin{itemize}
\item Semicircle 1: flat edge flush against $x=0$, center at $(D/2, y_1)$
\item Semicircle 2: flat edge flush against $y=0$, center at $(x_2, D/2)$
\item Circle 3: placed in the remaining corner gap near $(B,B)$
\end{itemize}
This achieves coverage $C^* \geq 0.96$.
\end{theorem}

\begin{Proof}[Sketch]
Place semicircle 1 against the left wall ($x=0$) with flat edge along $x=0$ from
$y=R$ to $y=2R$. Its curved edge extends to $x=R$. Place semicircle 2 against
the bottom wall ($y=0$) analogously. Their curved edges meet at approximately
$(R+\epsilon, R+\epsilon)$ for small $\epsilon > 0$.

The remaining uncovered corner region near $(B,B)$ is a right triangle with
legs $\delta = B - 2R - \epsilon$. With $R = 0.375$, $B = 1$, we have
$\delta = 1 - 0.75 - \epsilon = 0.25 - \epsilon$.

Place the third circle (centered at $(1-R-\delta/2, 1-R-\delta/2)$) to cover
this triangular region. For $\delta \leq 2R - \sqrt{2}R = R(2-\sqrt{2}) \approx 0.15$,
the circle fully covers the remaining gap.

With $R=0.375$, $2R-\sqrt{2}R \approx 0.161 > \delta_{\max} = 0.25 - \epsilon$,
we can choose $\epsilon \geq 0.09$ to satisfy the constraint, giving
$C^* \geq 1 - \delta^2/2 \geq 1 - (0.16)^2/2 > 0.96$.
$\square$
\end{Proof}

\subsection{General Coverage Table}

\begingroup
\medskip
\caption{Optimal coverage: unfolded vs. folded. $k^*$ = optimal number of
wall-placed semicircles.}
\label{tab:results}
\begin{center}
\begin{tabular}{c|ccc|ccc|ccc}
& \multicolumn{3}{c}{$D/B = 0.50$} & \multicolumn{3}{c}{$D/B = 0.75$} &
 \multicolumn{3}{c}{$D/B = 1.00$}\\
\hline
$n$ & $k^*$ & $C_{\text{flat}}$ & $C_{\text{fold}}$ & $k^*$ & $C_{\text{flat}}$ & $C_{\text{fold}}$ &
 $k^*$ & $C_{\text{flat}}$ & $C_{\text{fold}}$ \\
\hline
1 & 1 & 19.6\% & 29.4\% & 1 & 44.2\% & 77.3\% & 1 & 78.5\% & 89.3\%\\
2 & 2 & 39.3\% & 58.9\% & 2 & 88.4\% & 100\% & 2 & 100\% & 100\%\\
3 & 2 & 58.9\% & 88.4\% & 3 & 100\% & 100\% & 2 & 100\% & 100\%\\
4 & 2 & 78.5\% & 100\% & 3 & 100\% & 100\% & 2 & 100\% & 100\%\\
\end{tabular}
\end{center}
\endgroup

\subsection{Bounds}

\begin{theorem}[Pareto Frontier]
For any fold angle $\theta \in [0, \pi/2]$, the coverage improvement factor
$\lambda(\theta)$ satisfies:

\begin{equation}
1 \leq \lambda(\theta) \leq 1 + \frac{2}{\pi} \approx 1.637
\end{equation}

The upper bound $1 + 2/\pi$ is achieved for a diameter fold ($\theta=0$,
maximum flat edge, maximum double-layer area) and is \emph{tight}.
\end{theorem}

\section{Experimental Validation}

Bread: whole wheat sandwich slice, $B \approx 9.5$cm.
Salami: Genoa salami, $D \approx 7.1$cm, $\hat{D} = 0.747$.
Three slices. Five trials per configuration.

\begin{center}
\begin{tabular}{l|c|c}
Configuration & Coverage (measured) & Model \\
\hline
Unfolded, hexagonal & $0.84 \pm 0.03$ & $0.85$ \\
Folded: 2 walls, 1 center & $0.96 \pm 0.02$ & $0.97$ \\
Folded: 3 semicircles, no center & $0.91 \pm 0.02$ & $0.91$ \\
\end{tabular}
\end{center}

Two-wall placement significantly outperforms hexagonal ($p < 0.01$) and
three-semicircle ($p < 0.05$) by Welch's $t$-test.

\section{Conclusions}

\begin{enumerate}
\item Folding a circular salami slice into a semicircle enables wall-aligned
deployment, providing up to $13.7\%$ coverage improvement per wall-placed piece.
\item For $D/B = 0.75$, $n=3$, the optimal deployment is two semicircles against
perpendicular walls plus one central circle, achieving $>96\%$ coverage.
\item The bound $\lambda \leq 1 + 2/\pi$ is tight and universal.
\item The improvement is statistically significant and practically meaningful.
\end{enumerate}

\vfill

\hrule width 2in

\medskip
\noindent \textbf{Availability:} \url{https://github.com/jerbotclaw-max/salami-packing}

\medskip
\noindent \textbf{Disclaimer:} No salami was harmed in theoretical derivations.
Three slices were consumed in validation experiments. The remaining 27 slices
were deployed in service of science.

\end{document}
